\documentclass[12pt,nonacm,acmlarge,titlepage,screen]{acmart}

% make ACM settings and load extra packages
% settings specific to the ACM class

%% DO NOT MODIFY THIS FILE UNLESS IT IS FOR SUBMITTING TO AN ACM CONFERENCE/JOURNAL

%% Rights management information.  This information is sent to you
%% when you complete the rights form.  These commands have SAMPLE
%% values in them; it is your responsibility as an author to replace
%% the commands and values with those provided to you when you
%% complete the rights form.
% \setcopyright{acmlicensed}
\setcopyright{none}
% \copyrightyear{2018}
% \acmYear{2018}
\acmDOI{XXXXXXX.XXXXXXX}

%% These commands are for a JOURNAL article.
% \acmJournal{POMACS}
% \acmVolume{37}
% \acmNumber{4}
% \acmArticle{111}
% \acmMonth{8}

%% Submission ID.
%% Use this when submitting an article to a sponsored event. You'll
%% receive a unique submission ID from the organizers
%% of the event, and this ID should be used as the parameter to this command.
% \acmSubmissionID{123-A56-BU3}


%% citation style
% \citestyle{acmauthoryear}
\citestyle{acmnumeric}


%% ACM dates
% \received{\docDate}
% \received[revised]{12 March 2009}
% \received[accepted]{5 June 2009}
%% This file contains extra packages to use on top of the ACM template

%% STYLING ToC
\usepackage{titletoc}

%% SUB-FIGURES
\usepackage{subcaption}

%% ALGORITHMS
\usepackage{algorithm}
\usepackage{algorithmicx}
\usepackage{algpseudocode}

%% MORE EQUATIONS TOOLS
\usepackage{amsmath,mathtools}

%% CODE LISTINGS
\usepackage{listings}

% EASIER CROSS-REFERENCES
\usepackage[capitalise,nameinlink,noabbrev]{cleveref}

% TO GET ORDINAL NUMBERS
\usepackage[super]{nth}

% PAGE-BREAKABLE TABLES
\usepackage{longtable}

% IN-TABLE TEXT ALIGNMENTS
\usepackage{array}

% document details
% replace those with your details
\def\docTitle{Skin Lesion Classification on FPGA}
\def\docAuthor{Lewis \textsc{Turner}}
\def\docAuthorID{H00430662}
\def\docSupervisor{Dr. Rob \textsc{Stewart}, Ms Tess \textsc{Watt}}
\def\docDegree{BSc (Hons) Computer Science}
%\def\docDegree{MSc Cookery}
\def\docType{Research Report}
% \def\docType{Honours Dissertation}
% \def\docType{Masters Dissertation}
\def\docCourse{F20PA/B/C - Honours Project}
% \def\docCourse{F21RP - Research Methods and Project Planning}
% \def\docCourse{F21MP - Masters Project and Dissertation}

% load extra formatting
% formatting and package settings


% document details
\def\docDate{\nth{\the\day}~\ifcase\the\month\or
  January\or February\or March\or April\or May\or June\or
  July\or August\or September\or October\or November\or
  December\fi~\the\year}
\def\docInstitution{Heriot-Watt University}
\def\docSchool{School of Mathematical and Computer Sciences}
\author{\docAuthor}
\title{\docTitle}
\date{\docDate}

% change geometry of ACM template
\geometry{twoside=true, head=13pt, foot=2pc,
     paperwidth=21cm, paperheight=29.7cm,
     includeheadfoot,
     top=78pt, bottom=114pt, inner=81pt, outer=81pt,
     marginparwidth=4pc,heightrounded
 }%

% to make sections start on new pages
% \AddToHook{cmd/section/before}{\newpage}

% colors
% main text
\def\textColor{\color[HTML]{212121}} 
\def\secTextColor{\color[HTML]{424242}}
% footer and header
\def\headerColor{\color[HTML]{565656}} 
% title, chapters, sections
\def\titleColor{\color[HTML]{212121}}

% set line spacing
% \onehalfspacing
% \singlespacing

% set table cells padding
\def\arraystretch{1.5}

% set headings format
% unset to follow ACM format
% \titleformat{\chapter}{\singlespacing\normalfont\huge\bfseries\titleColor}{\chaptertitlename\ \thechapter}{1em}{}
% \titleformat{\chapter}{\singlespacing\normalfont\huge\bfseries\titleColor}{\thechapter}{1em}{}
% \titleformat{\section}{\singlespacing\normalfont\large\bfseries\titleColor}{\thesection}{1em}{}
% \titleformat{\subsection}{\singlespacing\normalfont\large\bfseries}{\thesubsection}{1em}{}
% \titleformat{\subsubsection}{\singlespacing\normalfont\normalsize\bfseries}{\thesubsubsection}{1em}{}
% \titlespacing*{\chapter}{0pt}{0pt}{10pt}

% set paragraph indent
% unset to follow ACM format
% \setlength{\parindent}{18pt} 

% set the text colour
% \textColor 

% set hyperlinks + pdf metadata
\hypersetup{
    pdftitle={\docTitle},
    pdfauthor={\docAuthor},
    pdfkeywords={computer science dissertation},
    colorlinks=true,%
    citecolor=[HTML]{152845},%
    linkcolor=[HTML]{152845},%
    urlcolor=[HTML]{1976D2}
}

% styling the code listings
\definecolor{codegreen}{rgb}{0,0.6,0}
\definecolor{codegray}{rgb}{0.5,0.5,0.5}
\definecolor{codepurple}{rgb}{0.58,0,0.82}
\definecolor{backcolour}{rgb}{0.95,0.95,0.92}

\lstdefinestyle{mystyle}{
    backgroundcolor=\color{backcolour},   
    commentstyle=\color{codegreen},
    keywordstyle=\color{magenta},
    numberstyle=\tiny\color{codegray},
    stringstyle=\color{codepurple},
    basicstyle=\ttfamily\singlespacing,
    breakatwhitespace=false,         
    breaklines=true,                 
    captionpos=b,                    
    keepspaces=true,                 
    numbers=left,                    
    numbersep=5pt,                  
    showspaces=false,                
    showstringspaces=false,
    showtabs=false,                  
    tabsize=2
}
\lstset{style=mystyle}

% reset names for code listing
\renewcommand{\lstlistingname}{Code Listing}
\renewcommand{\lstlistlistingname}{List of \lstlistingname s}

% set cref for listing
\crefname{lstlisting}{\lstlistingname}{\lstlistingname s}
\Crefname{lstlisting}{\lstlistingname}{\lstlistingname s}

% set cross-references labels with cleveref
\crefname{appendixchapter}{Appendix}{Appendices}
\Crefname{appendixchapter}{Appendix}{Appendices}


% set custom headers/footers
% unset to follow ACM template
% \pagestyle{fancy}
% \fancyhf{}
% \renewcommand{\headrulewidth}{0pt}
% \renewcommand{\footrulewidth}{0pt}
% \fancyhead[R]{\headerColor \docType}
% \fancyhead[L]{\headerColor \docAuthor}
% \fancyfoot[C]{\headerColor \thepage}
% % Redefine the plain page style for chapter pages
% \fancypagestyle{plain}{%
%     \fancyhf{}%
%     \fancyfoot[C]{\headerColor \thepage}
% }

% custom section heading for front-matter content with addition to the ToC
\newcommand{\frontmattersection}[2]{
    \thispagestyle{plain} % remove the header on page
    % \vspace*{1cm} % to align vertically
    \phantomsection
    \begin{center} 
        \textsc{#1} 
        \addcontentsline{toc}{section}{#1}
    \end{center}
    #2 % Main Text
    \clearpage % leave blank space underneath
}

\renewcommand{\abstract}[1]{
    \frontmattersection{Abstract}{#1}
}

\newcommand{\acknowledgements}[1]{
    \frontmattersection{Acknowledgements}{#1}
}

% formatting the table of content
\dottedcontents{section}[1em]{\bfseries}{1.2em}{0pc}
\dottedcontents{subsection}[3em]{}{2.2em}{1pc}

\renewcommand{\contentsname}{Table of Contents}

\begin{document}

%%%%%%%%%%%%%%%%%%%%%%%%%%%%%%%%%%%%%%%%%%%%%%%%%%
% FRONT MATTER
%%%%%%%%%%%%%%%%%%%%%%%%%%%%%%%%%%%%%%%%%%%%%%%%%%
\begin{titlepage}

\begin{center}
	
	% Title (and subtitle)
	~\\[0.2cm]
	{\huge \titleColor \textbf{\docTitle}}\\
	~\\[0.7cm]
	
	% Author(s)
	{\Large \textit{by} \docAuthor}\\[0.7cm]
	{\large \docType}\\[1cm]
	{\large \textit{Submitted in partial fulfilment for the degree of}\\[0.2cm] \docDegree}\\[0.7cm]
	{\large \textit{Supervised by} {\Large \docSupervisor}}
	~\\[1cm]
	
	% Organisation	
	% ~\\[0cm]
        % logo / coat of arms
        \begin{figure}[h]
            \centering
		\includegraphics[height=3cm,keepaspectratio]{images/hwcoa_bw.jpg}
        \Description{Heriot-Watt University's coat-of-arms}
	\end{figure}
	\begin{large}
		\textsc{\docInstitution}\\
		\docSchool\\
% 		Department\\
	\end{large}
	~\\[1cm]
	
	% Date
	{\large \docDate}\\
	~\\[1.5cm]
	
	% Copyright
	\begin{minipage}{0.7\textwidth}
	% \singlespacing
	\begin{footnotesize}
	The copyright in this dissertation is owned by the author. Any quotation from the dissertation or use of any of the information contained in it must acknowledge it as the source of the quotation or information.
	\end{footnotesize}
	\end{minipage}
	
\end{center}

\end{titlepage}

\newpage
\pagenumbering{roman}

% authorship declaration
% this declaration is automatically filled with the details provided in main.tex

\thispagestyle{plain} % remove the header on page

\phantomsection
\begin{center} 
    \textsc{Declaration of Authorship} 
    \addcontentsline{toc}{section}{Declaration of Authorship}
\end{center}


\begin{center} 
    \begin{tabular}{|p{0.3\textwidth}|p{0.675\textwidth}|}
        \hline
        \textbf{Course code and name} & \docCourse \\
        \hline
        \textbf{Type of assessment} & Individual \\
        \hline
        \textbf{Document title} & \docTitle \\
        \hline
        \textbf{Student name} & \docAuthor \\
        \hline
        \textbf{Student ID} & \docAuthorID \\
        \hline
    \end{tabular}
\end{center}

\bigskip

\begin{longtable}{|>{\raggedright\arraybackslash}p{\textwidth}|}
    \hline
    \endfirsthead
    \hline
    \multicolumn{1}{|r|}{\footnotesize\it (form continued from previous page)} \\
    \endhead
    \multicolumn{1}{|r|}{\footnotesize\it (form continues on next page)} \\
    \hline
    \endfoot
    \hline
    \endlastfoot
    
        \textbf{Statement on Generative AI Use:}
        \\
        % EDIT this statement to indicate what use you made of Generative AI for your project work and in preparation of this document. If your project is about developing or experimenting with Generative AI, you still need to indicate your Generative AI use beyond the developments and experiments. If you have not used Generative AI, indicate it explicitly.
        % \\
        
        \textbullet~ Google Gemini 3.8 Flash (Extended Thinking) to assist in searching for relevant research papers.
        \\
        \textbullet~ Google Gemini 4.0 Argon (thinking: Low - High) to assist setting up FPGA for use over Wifi.
        \\
    \hline
    \multicolumn{1}{l}{}
    \\
    \hline
        \textbf{By signing this form:}
        \\
        
        \textbullet~ I declare that the work I have submitted for this assessment is my own original work. I have not presented the ideas, words, data, code, designs, analyses, or other intellectual contributions of another person or system as my own. 
        \\
        
        \textbullet~ My submission is expressed in my own words except where the work of others has been appropriately acknowledged, cited, and referenced. All sources used in the preparation of this assessment, including books, journal articles, websites, software tools, and Generative AI tools, have been properly acknowledged.
        \\
        
        \textbullet~ Where Generative AI tools have been used, I have complied with the course’s stated policies regarding Generative AI use. I have clearly and accurately disclosed in the above statement box all uses of Generative AI in the assessment, including the specific tools used, the purpose for which they were used, and the extent to which any generated content influenced the final submission.
        \\
        
        \textbullet~ I understand that I remain fully responsible for the content, accuracy, integrity, and originality of my submission, regardless of any assistance received from Generative AI or other tools.
        \\

        \textbullet~ I confirm that I have read, understood and followed the University's Regulations on academic misconduct as published on the \href{https://www.hw.ac.uk/uk/services/academic-registry/academic-integrity/academic-misconduct.htm}{University's website}, and that I am aware of the penalties that I will face should I not adhere to the University Regulations.
        \\
        
        \textbullet~ I confirm that I have read, understood and avoided the different types of plagiarism explained in the University guidance on \href{https://heriotwatt.sharepoint.com/sites/skillshub/SitePages/Academic-Integrity-and-Plagiarism.aspx}{Academic Integrity and Plagiarism}.
        \\
        
        \\

        \noindent \textbf{Student signature}: \docAuthor
        \\
        
        \\
        
        \noindent \textbf{Date}: \docDate
        \\
    \hline
\end{longtable}

% \noindent \textbf{Your work will not be marked if a signed copy of this form is not included with your submission.}



\clearpage

% Acknowledgements
\acknowledgements{
Acknowledgements text.
}

% Abstract
\abstract{
Abstract text.
}

% table of contents

\newpage

\thispagestyle{plain} % remove the header on page

\setcounter{tocdepth}{2} % set depth of ToC, 0 = cha, 1 = sec, etc.

\tableofcontents % the actual table of contents command

\clearpage

	
%%%%%%%%%%%%%%%%%%%%%%%%%%%%%%%%%%%%%%%%%%%%%%%%%%
% MAIN BODY
%%%%%%%%%%%%%%%%%%%%%%%%%%%%%%%%%%%%%%%%%%%%%%%%%%
\newpage
\pagenumbering{arabic}

% Introduction
%%%%%%%%%%%%%%%%%%%%
\section{Introduction} \label{cha:intro}
Summary of project motivation \& context

    \subsection{Background and Motivation} \label{sec:intro_motiv}
    \textbullet~ Using papers (mainly) from literature review to provide background context\\
    \textbullet~ Motivations for doing this project\\
    \textbullet~ Identifying problems to be solved to achieve objectives

    \subsection{Aim and Objectives} \label{sec:intro_aim}
    State project aims\\
    Table of objectives with IDs to reference later

% Plan of the document.

% Background Research
%%%%%%%%%%%%%%%%%%%%
\section{Literature Review} \label{cha:lit_rev}
Grouping topic summaries, critical assessment, gaps\\
Possibly use RADAR?

    % How I searched for papers, what tools did I use and how did I verify my sources
    \subsection{Search Methods}
    I used Google Gemini 3.8 Flash on thinking mode High (October 2026 version) to source the most
    recent academic papers on a particular topic. I would ask it to create a table consisting of the
    papers: title, year, summary, DOI or similar URL. This allowed me to leverage the LLM's speed for
    searching the internet like a more powerful keyword search.
    \\
    \\
    I occasionally used Google Scholar as well to search a specific paper title provided by Gemini to
    find the original source PDF/website.

    % what to do:
    % Identify key topics, themes, points
    %
    % For each summary of literature \& critical assessment of topic containing:\\
    %     \textbullet~ How already published research fits together\\
    %     \textbullet~ Areas of consensus \& disagreement\\
    %     \textbullet~ Areas of controversy/where further research is required


    % RADAR: Relevance, Authority, Date, Appearance, Reason for writing

    % Papers on recent skin classification models
    \subsection{Skin Lesion Classification Models}
    Many different types of ML models trained to identify skin lesions from images:
    \\
    Multi-modal approach on an FPGA which combines image features with patient metadata \cite{jesmdt_2025}.
    Uses MobileNetV2 architecture \cite{mobilenet_2017}.
    \\


    % Papers more generally on methodology for putting ML models on FPGA
    \subsection{Machine Learning Models on FPGA}
    Very popular fundamental paper presenting MobileNet architecture \cite{mobilenet_2017} % find MobileNetV2 paper?
    
    \subsection{Compressing Machine Learning Models to FPGA}
    Papers on useful compression techniques for certain model types
    
    \subsection{Further Research Recommendations}
    Identify gaps, need to repeat any research

    % Sources for literature/statements
    \subsection{Acknowledgements}
    Skin lesion paper 1 \cite{jesmdt_2025}
    % \\
    % Example citation to stop LaTeX complaining: \cite{turing_computing_2009}. % example citation

\section{Methodology} \label{sec:method}
Detailed description of necessary methods, requirements, research questions that will be employed/
needed to complete the project

    \subsection{Requirements}
    MOSCOW\\
    Enumerated to be easier to reference again

    \subsection{Plan for Execution}
    Explain how I plan to use experiments to fulfil all the requirements \& objectives

    \subsection{Experiment Planning}
    How I plan to conduct experiment(s) and what research question does it address/answer (detailed)

\section{Feasibility} \label{sec:feasiblity}

    \subsection{Preliminary Work}
    Summary of work carried out to prepare for this project to validate feasibility.

    \subsection{Evaluation Strategy}
    Identify for the project:
    \\
        \textbullet~ What needs evaluating
        \\
        \textbullet~ Why it should be evaluated
        \\
        \textbullet~ How it should be evaluated
        \\

\section{Conclusions}
Short summary of this doc

% Rest of the document...
%%%%%%%%%%%%%%%%%%%%


%%%%%%%%%%%%%%%%%%%%%%%%%%%%%%%%%%%%%%%%%%%%%%%%%%
% BACK MATTER
%%%%%%%%%%%%%%%%%%%%%%%%%%%%%%%%%%%%%%%%%%%%%%%%%%

% References
% list of references

\newpage
\renewcommand{\bibname}{References} % change name to References
% \phantomsection
% \addcontentsline{toc}{section}{References}
\bibliography{references}
% \bibliographystyle{plainnat}

% using Harvard style
% \bibliographystyle{agsm}

% using ACM style (defined in acm_settings.tex)
\bibliographystyle{ACM-Reference-Format}

% using APA style (requires apacite package)
% \bibliographystyle{apacite}

\section{Project Management}

    \subsection{Gantt Chart}
    Plan of activities to be carried out in semester 1, then future plan for semester 2

    \subsection{Analysis of Risks}
    Including mitigation plans

    \subsection{PLES}
    A well-researched consideration of any Professional, Legal, Ethical, and Social Issues
    pertinent to the project. (e.g. codes of conduct (BCS), codes of practice, standards,
    computer law, ethical decision making, intellectual property, social aspects, copyright,
    data protection, and so on).

\section{Gen AI Summary}
Summarising use of GenAI in project/deliverables

% Appendices
\newpage
\appendix

\crefalias{chapter}{appendixchapter} % for appendix crossreferences

% Rest of appendices...

\end{document}
